%% ================================================================
%% RAPPORT PFE — Mohamed Yassine Marzouk
%% Association El Awael — Master Leadership Managerial 4.0
%% Compiler : XeLaTeX (deux passes)
%% ================================================================
\documentclass[12pt,a4paper,openany]{report}

%% ── Langue francaise ───────────────────────────────────────────
\usepackage[french]{babel}

%% ── Police Times New Roman (XeLaTeX) ───────────────────────────
\usepackage{fontspec}
\setmainfont{Times New Roman}
\setsansfont{Arial}
\setmonofont{Courier New}

%% ── Mise en page ───────────────────────────────────────────────
\usepackage[top=2.5cm,bottom=2.5cm,left=3cm,right=2.5cm,headheight=26pt]{geometry}

%% ── Interligne 1,5 ─────────────────────────────────────────────
\usepackage{setspace}
\onehalfspacing

%% ── Couleurs ───────────────────────────────────────────────────
\usepackage[dvipsnames]{xcolor}
\definecolor{bleu}{RGB}{27,58,107}
\definecolor{bleuClair}{RGB}{210,222,240}
\definecolor{grisTxt}{RGB}{80,80,80}
\definecolor{vertOK}{RGB}{34,120,34}
\definecolor{rougeKO}{RGB}{170,30,30}

%% ── En-tetes et pieds de page ──────────────────────────────────
\usepackage{fancyhdr}
\pagestyle{fancy}
\fancyhf{}
\fancyhead[L]{\small\color{grisTxt}\itshape\nouppercase{\leftmark}}
\fancyhead[R]{\small\color{grisTxt}\itshape M.Y. Marzouk}
\fancyfoot[C]{\small\thepage}
\renewcommand{\headrulewidth}{0.4pt}
\renewcommand{\chaptermark}[1]{\markboth{#1}{}}

%% ── Titres ─────────────────────────────────────────────────────
\usepackage{titlesec}
\titleformat{\chapter}[display]
  {\normalfont\fontsize{14}{18}\bfseries\color{bleu}}
  {\fontsize{13}{16}\bfseries\chaptertitlename~\thechapter}{8pt}
  {\fontsize{14}{18}\bfseries}
\titlespacing{\chapter}{0pt}{0pt}{22pt}

\titleformat{\section}
  {\normalfont\fontsize{14}{17}\bfseries\color{bleu}}
  {\thesection}{10pt}{}
\titlespacing{\section}{0pt}{14pt}{6pt}

\titleformat{\subsection}
  {\normalfont\fontsize{12}{15}\bfseries\color{bleu!80!black}}
  {\thesubsection}{8pt}{}
\titlespacing{\subsection}{0pt}{10pt}{4pt}

\titleformat{\subsubsection}
  {\normalfont\fontsize{12}{15}\itshape\color{grisTxt}}
  {\thesubsubsection}{6pt}{}

%% ── Figures et tableaux ────────────────────────────────────────
\usepackage{graphicx}
\usepackage{float}
\usepackage{caption}
\captionsetup{labelfont={bf,color=bleu},textfont=it,justification=centering,skip=5pt}
\usepackage{booktabs}
\usepackage{array}
\usepackage{tabularx}
\usepackage{multirow}
\usepackage{longtable}
\usepackage{chngcntr}
\counterwithin{figure}{chapter}
\counterwithin{table}{chapter}

%% ── Boites colorees (tcolorbox) ────────────────────────────────
\usepackage{tcolorbox}
\tcbuselibrary{skins,breakable}

\newtcolorbox{encadre}[2][]{
  colback=bleuClair!30,colframe=bleu,boxrule=1pt,arc=3pt,
  title=#2,fonttitle=\bfseries\color{bleu},#1}

\newtcolorbox{boitecitation}{
  colback=gray!8,colframe=gray!40,boxrule=0.5pt,arc=2pt,
  left=8pt,right=8pt}

%% ── Listes ─────────────────────────────────────────────────────
\usepackage{enumitem}
\setlist[itemize]{itemsep=2pt,topsep=4pt}
\setlist[enumerate]{itemsep=2pt,topsep=4pt}

%% ── Notes de bas de page ───────────────────────────────────────
\usepackage[bottom]{footmisc}
\renewcommand{\footnotesize}{\fontsize{10}{12}\selectfont}

%% ── Symboles ───────────────────────────────────────────────────
\usepackage{pifont}
\newcommand{\cmark}{\ding{51}}
\newcommand{\xmark}{\ding{55}}

%% ── Hyperliens ─────────────────────────────────────────────────
\usepackage[colorlinks=true,linkcolor=bleu,citecolor=bleu,urlcolor=bleu,
  pdftitle={Rapport PFE Mohamed Yassine Marzouk},
  pdfauthor={Mohamed Yassine Marzouk}]{hyperref}

%% ── Image securisee (placeholder si fichier absent) ────────────
\newcommand{\safeimg}[3][0.9\textwidth]{%
  \IfFileExists{#2}{\includegraphics[width=#1]{#2}}%
  {\fbox{\parbox{#1}{\centering\color{grisTxt}\itshape[Image~:~#3]}}}}

%% \og et \fg sont deja definis par babel French — pas besoin de les redefinir

%% ================================================================
\begin{document}
\pagenumbering{gobble}

%% ================================================================
%%  PAGE DE GARDE
%% ================================================================
\begin{titlepage}
\begin{center}

  \IfFileExists{assets/logo/logo_elawael.png}{%
    \includegraphics[width=3.2cm]{assets/logo/logo_elawael.png}\\[0.4cm]}{}

  {\large\bfseries Universite Abdelmalek Essaadi}\\[0.1cm]
  {\large Faculte des Sciences Juridiques, Economiques et Sociales}\\[0.05cm]
  {\normalsize Tetouan}\\[0.5cm]

  \textcolor{bleu}{\rule{\textwidth}{2pt}}\\[0.2cm]
  {\large\bfseries\color{bleu} Filiere~: Master Leadership Managerial 4.0}\\[0.2cm]
  \textcolor{bleu}{\rule{\textwidth}{2pt}}

  \vspace{0.7cm}
  {\Large\bfseries RAPPORT DE STAGE DE FIN D'ETUDES}
  \vspace{0.6cm}

  \begin{tcolorbox}[colback=bleu!8,colframe=bleu,boxrule=1.5pt,arc=5pt,halign=center,width=\textwidth]
    \vspace{0.2cm}
    {\fontsize{13}{17}\bfseries\color{bleu}
      Innovation manageriale et transformation digitale\\
      dans les associations sociales :\\[0.3cm]
      Entre performance organisationnelle et impact humain}
    \vspace{0.2cm}
  \end{tcolorbox}

  \vspace{0.6cm}
  {\normalsize Realise au sein de~:}\\[0.1cm]
  {\large\bfseries Association El Awael --- Centre de Martil}\\[0.1cm]
  {\normalsize Periode~: 12 mars 2026 --- 12 juin 2026}

  \vspace{1cm}
  \begin{minipage}[t]{0.47\textwidth}
    \raggedright
    {\bfseries Realise par~:}\\[0.1cm]
    Mohamed Yassine Marzouk\\
    Master Leadership Managerial 4.0\\
    Annee universitaire 2025--2026
  \end{minipage}
  \hfill
  \begin{minipage}[t]{0.47\textwidth}
    \raggedleft
    {\bfseries Encadrant academique~:}\\[0.1cm]
    M. Abdelmonaim Tlidi\\[0.3cm]
    {\bfseries Responsable d'accueil~:}\\[0.1cm]
    Mme Kaoutar Dghoughi\\
    {\small Directrice, Centre de Martil}
  \end{minipage}

  \vfill
  \textcolor{bleu}{\rule{\textwidth}{1pt}}\\[0.2cm]
  {\large\bfseries Annee universitaire 2025--2026}
\end{center}
\end{titlepage}

%% ================================================================
%%  DEDICACE
%% ================================================================
\cleardoublepage
\thispagestyle{empty}
\vspace*{5cm}
\begin{center}
  \textcolor{bleu}{\rule{0.4\textwidth}{0.5pt}}\\[0.5cm]
  {\large\itshape\bfseries Dedicace}\\[0.5cm]
  \textcolor{bleu}{\rule{0.4\textwidth}{0.5pt}}
\end{center}
\vspace{1cm}
\begin{flushright}
\begin{minipage}{0.7\textwidth}
\itshape
A mes parents, source inepuisable d'amour et de soutien,\\
qui ont tout sacrifie pour que je puisse atteindre mes ambitions.\\[0.4cm]
A mes enseignants de la Faculte des Sciences Juridiques,\\
Economiques et Sociales de Tetouan,\\
qui ont forge mon regard sur le monde managerial.\\[0.4cm]
A chaque enfant en situation de handicap suivi par\\
l'Association El Awael, raison d'etre de ce travail.\\[0.5cm]
\raggedleft --- Mohamed Yassine Marzouk
\end{minipage}
\end{flushright}

%% ================================================================
%%  REMERCIEMENTS
%% ================================================================
\cleardoublepage
\thispagestyle{empty}
\begin{center}
{\large\bfseries\color{bleu} Remerciements}\\[0.2cm]
\textcolor{bleu}{\rule{0.5\textwidth}{0.5pt}}
\end{center}
\vspace{0.5cm}

Je tiens a exprimer ma profonde gratitude a toutes les personnes
qui ont contribue, de pres ou de loin, a la realisation de ce stage
et de ce rapport.

\vspace{0.3cm}
Mes remerciements les plus sinceres vont a \textbf{Monsieur Abdelmonaim Tlidi},
mon encadrant academique, pour ses precieux conseils, sa disponibilite
constante et son accompagnement rigoureux tout au long de ce travail.

\vspace{0.3cm}
Je remercie chaleureusement \textbf{Madame Kaoutar Dghoughi}, directrice du
centre de Martil de l'Association El Awael, pour m'avoir accueilli avec
bienveillance au sein de son equipe et pour la confiance qu'elle m'a accordee.

\vspace{0.3cm}
Je suis particulierement reconnaissant envers l'ensemble de
\textbf{l'equipe educative du centre de Martil} --- cadres encadrants,
educateurs specialises et personnel administratif --- pour leur accueil
chaleureux et leur engagement quotidien aupres des enfants en situation
de handicap.

\vspace{0.3cm}
Je remercie egalement le corps professoral de la \textbf{Faculte des Sciences
Juridiques, Economiques et Sociales} de l'Universite Abdelmalek Essaadi
de Tetouan, et plus particulierement les enseignants du Master Leadership
Managerial 4.0, pour les solides bases intellectuelles transmises.

\vspace{0.5cm}
\begin{flushright}
\itshape Mohamed Yassine Marzouk\\
Martil, juin 2026
\end{flushright}

%% ================================================================
%%  RESUME / ABSTRACT
%% ================================================================
\cleardoublepage
\thispagestyle{empty}
\begin{center}
{\large\bfseries\color{bleu} Resume}\\[0.2cm]
\textcolor{bleu}{\rule{0.5\textwidth}{0.5pt}}
\end{center}
\vspace{0.4cm}

Ce rapport de stage de fin d'etudes, realise au sein de l'Association El
Awael au centre de Martil (Maroc) du 12 mars au 12 juin 2026, explore
les conditions d'une transformation digitale au service de l'innovation
manageriale dans une association a vocation sociale. La mission a conduit
a la conception et au deploiement d'une \textbf{plateforme de gestion numerique
integree} --- ElAwael Digital Platform --- couvrant l'authentification
securisee, la gestion des dossiers des beneficiaires en situation de
handicap, le suivi pedagogique individualise (PEI), les tableaux de bord
decisionnels (KPI), la simulation strategique et le journal d'audit.

\vspace{0.3cm}
L'analyse theorique mobilise les concepts du Management 4.0, de
l'innovation manageriale et de la transformation digitale, appliques a la
specificite du secteur associatif marocain. Le diagnostic organisationnel
revele des dysfonctionnements critiques que la solution technique resout
de maniere structuree, avec des gains mesurables en performance et en
qualite du suivi des beneficiaires.

\vspace{0.3cm}
\textbf{Mots-cles~:} Transformation digitale, Innovation manageriale,
Management 4.0, Association sociale, PEI, KPI, Firebase,
Performance organisationnelle, Handicap.

\vspace{0.8cm}
\begin{center}
{\large\bfseries\color{bleu} Abstract}\\[0.2cm]
\textcolor{bleu}{\rule{0.5\textwidth}{0.5pt}}
\end{center}
\vspace{0.4cm}

This internship report, conducted at the El Awael Association's Martil
centre (Morocco) from March 12 to June 12, 2026, examines the conditions
for digital transformation serving managerial innovation in a social
welfare organization. The mission led to the design and deployment of an
integrated digital management platform covering secure authentication,
beneficiary file management, individualized educational planning (IEP)
monitoring, KPI dashboards, strategic simulation and audit logging.

\vspace{0.3cm}
\textbf{Keywords~:} Digital transformation, Managerial innovation,
Management 4.0, Social association, IEP, KPI, Firebase,
Organizational performance, Disability.

%% ================================================================
%%  LISTE DES ABREVIATIONS
%% ================================================================
\cleardoublepage
\thispagestyle{empty}
\begin{center}
{\large\bfseries\color{bleu} Liste des abreviations}\\[0.2cm]
\textcolor{bleu}{\rule{0.5\textwidth}{0.5pt}}
\end{center}
\vspace{0.4cm}

\begin{longtable}{@{}>{\bfseries}p{3.5cm}p{11cm}@{}}
\toprule
\textbf{Abreviation} & \textbf{Signification}\\
\midrule
ABLLS & Assessment of Basic Language and Learning Skills\\
API   & Application Programming Interface\\
CIN   & Carte d'Identite Nationale\\
CRUD  & Create, Read, Update, Delete\\
FSJES & Faculte des Sciences Juridiques, Economiques et Sociales\\
IA    & Intelligence Artificielle\\
JS    & JavaScript\\
KPI   & Key Performance Indicator (Indicateur Cle de Performance)\\
ONG   & Organisation Non Gouvernementale\\
PEI   & Projet Educatif Individualise\\
RTL   & Right-To-Left (ecriture droite a gauche)\\
SI    & Systeme d'Information\\
SOG   & Score Organisationnel Global\\
TIC   & Technologies de l'Information et de la Communication\\
UX/UI & User Experience / User Interface\\
\bottomrule
\end{longtable}

%% ================================================================
%%  TABLES
%% ================================================================
\cleardoublepage
\pagenumbering{roman}
\setcounter{page}{1}
\tableofcontents

\cleardoublepage
\listoffigures
\addcontentsline{toc}{chapter}{Liste des figures}

\cleardoublepage
\listoftables
\addcontentsline{toc}{chapter}{Liste des tableaux}

%% ================================================================
%%  CORPS
%% ================================================================
\cleardoublepage
\pagenumbering{arabic}
\setcounter{page}{1}

%% ================================================================
%%  INTRODUCTION GENERALE
%% ================================================================
\chapter*{Introduction generale}
\addcontentsline{toc}{chapter}{Introduction generale}
\markboth{Introduction generale}{Introduction generale}

\begin{boitecitation}
\itshape
\og{}L'innovation est l'instrument specifique de l'entrepreneuriat,
l'acte qui confere aux ressources une nouvelle capacite a creer de la
richesse.\fg{}\\
\raggedleft --- Peter Drucker, 1985
\end{boitecitation}

\vspace{0.5cm}
Dans un monde en mutation acceleree, marque par l'omniprésence du
numerique et l'exigence croissante de transparence et d'efficacite,
aucun secteur n'echappe a la necessite de se reinventer. Le secteur
associatif marocain, traditionnellement ancre dans des pratiques de
solidarite et d'accompagnement humain, se trouve aujourd'hui au
carrefour d'une double obligation~: preserver sa mission sociale
fondamentale tout en integrant les apports transformateurs des
technologies de l'information et de la communication.

\vspace{0.3cm}
L'Association El Awael, implantee dans la region de Tetouan avec des
centres a Martil, M'diq et Fnideq, incarne cette realite avec une
acuite particuliere. Dediee a l'accompagnement des enfants en situation
de handicap et de leurs familles, l'association dispose d'un capital
humain engage et d'une mission d'une importance capitale pour les
populations les plus vulnerables. Pourtant, elle se heurtait, avant ce
stage, a des limites structurelles~: une gestion documentaire encore
largement manuelle, une tracabilite insuffisante, l'absence d'outils
decisionnels et une securite des donnees preoccupante.

\vspace{0.3cm}
C'est dans ce contexte que s'inscrit le present stage de fin d'etudes,
realise au centre de Martil du 12 mars au 12 juin 2026, dans le cadre
du Master Leadership Managerial 4.0 de la Faculte des Sciences
Juridiques, Economiques et Sociales de l'Universite Abdelmalek Essaadi
de Tetouan.

\begin{encadre}{Problematique centrale}
\textbf{Comment la transformation digitale, articulee aux principes du
Management 4.0, peut-elle constituer un levier d'innovation manageriale
pour une association sociale, en conciliant performance organisationnelle
et preservation de l'impact humain aupres des beneficiaires~?}
\end{encadre}

\vspace{0.3cm}
Pour repondre a cette question, ce rapport est structure en
\textbf{trois parties}~:

\begin{description}
  \item[Partie I] Presentation de l'Association El Awael et cadre
    theorique de la transformation digitale (Chapitres 1 et 2).
  \item[Partie II] Diagnostic organisationnel et conception de la
    solution digitale (Chapitres 3 et 4).
  \item[Partie III] Realisation de la plateforme et evaluation de
    l'impact (Chapitres 5 et 6).
\end{description}

%% ================================================================
\part{Cadre general et fondements theoriques}

%% ================================================================
\chapter{Presentation de l'Association El Awael et contexte du stage}

\section{Le secteur associatif marocain}

Le tissu associatif marocain a connu une croissance exponentielle au
cours des trois dernieres decennies. Le Maroc compte aujourd'hui plus
de 230\,000 associations legalement enregistrees\footnote{Source~:
Secretariat General du Gouvernement, Rapport annuel sur le mouvement
associatif, 2024.}, couvrant des champs aussi varies que l'education,
la sante, le developpement rural, la culture et l'aide aux personnes
en situation de vulnerabilite.

\vspace{0.3cm}
La majorite des associations a vocation sociale souffrent de contraintes
communes~: dependance aux financements publics, turn-over important des
ressources humaines benevoles, absence de systemes d'information
formalises, et difficultes a rendre compte de leur impact de maniere
rigoureuse aupres des partenaires institutionnels.

\begin{table}[H]
\centering
\caption{Caracteristiques du secteur associatif marocain}
\label{tab:secteur}
\begin{tabularx}{\textwidth}{>{\bfseries}lX}
\toprule
Indicateur & Donnees\\
\midrule
Associations enregistrees & $>$ 230\,000 (2024)\\
Associations actives estimees & 45\,000 a 60\,000\\
Domaine social et caritatif & $\approx$ 35\% du total\\
Associations sans systeme d'info. & $>$ 70\% (estimation)\\
\bottomrule
\end{tabularx}
\vspace{0.1cm}
\small\textit{Source~: SGG, INDH, compilation auteur, 2024--2026.}
\end{table}

\section{Identite et historique de l'Association El Awael}

\subsection{Fondation et mission}

L'Association El Awael est une organisation a but non lucratif marocaine
dediee a la \textbf{prise en charge, a l'education specialisee et a la
rehabilitation des enfants en situation de handicap}. Sa mission
fondamentale s'articule autour de trois axes~:

\begin{enumerate}
  \item \textbf{L'inclusion educative}~: accompagner chaque beneficiaire
    vers le developpement de ses capacites cognitives, langagieres et
    motrices, a travers des programmes individualises adaptes a chaque
    profil de handicap~;
  \item \textbf{Le soutien aux familles}~: offrir un espace d'ecoute,
    de formation et d'accompagnement aux parents et tuteurs~;
  \item \textbf{Le plaidoyer social}~: sensibiliser la communaute et
    les institutions aux droits des personnes en situation de handicap.
\end{enumerate}

\subsection{Implantation geographique}

L'Association El Awael dispose de \textbf{trois centres} dans la region
de M'diq-Fnideq~:

\begin{table}[H]
\centering
\caption{Les trois centres de l'Association El Awael}
\label{tab:centres}
\begin{tabularx}{\textwidth}{lXXl}
\toprule
\textbf{Centre} & \textbf{Role} & \textbf{Specificites} & \textbf{Statut}\\
\midrule
Fnideq & Siege social & Administration centrale & Siege principal\\
\textbf{Martil} & Centre operationnel & Suivi PEI, presences & \textbf{Centre du stage}\\
M'diq & Centre complementaire & Rehabilitation, ateliers & Antenne regionale\\
\bottomrule
\end{tabularx}
\vspace{0.1cm}
\small\textit{Source~: Documentation interne de l'Association El Awael, 2026.}
\end{table}

\subsection{Structure organisationnelle}

\begin{figure}[H]
\centering
\begin{tcolorbox}[colback=bleuClair!20,colframe=bleu,boxrule=0.8pt,arc=3pt,width=0.9\textwidth]
\centering\small
{\bfseries\color{bleu} Bureau Executif (Fnideq)}\\[0.25cm]
$\downarrow$\\[0.15cm]
{\bfseries Directrice du Centre de Martil}\\
{\small Mme Kaoutar Dghoughi}\\[0.25cm]
$\downarrow$\\[0.15cm]
\begin{tabular}{ccc}
\fbox{\parbox{3cm}{\centering\small Equipe pedagogique\\(Educateurs specialises)}} & \quad &
\fbox{\parbox{3cm}{\centering\small Personnel\\administratif}}
\end{tabular}\\[0.25cm]
$\downarrow$\\[0.15cm]
{\bfseries Beneficiaires et familles}
\end{tcolorbox}
\caption{Organigramme simplifie du centre de Martil}
\label{fig:organigramme}
\small\textit{Source~: Elaboration de l'auteur d'apres observations, 2026.}
\end{figure}

\section{Les beneficiaires et le Projet Educatif Individualise (PEI)}

Le centre de Martil accueille des enfants et adolescents presentant
des profils varies de handicap~: deficiences intellectuelles, troubles
du spectre autistique (TSA), handicaps moteurs et troubles sensoriels.

\vspace{0.3cm}
L'outil pedagogique central est le \textbf{Projet Educatif Individualise
(PEI)}, document etabli pour chaque beneficiaire, definissant~:

\begin{encadre}{Composantes d'un PEI}
\begin{itemize}
  \item Les objectifs a court, moyen et long terme (cognitifs, moteurs,
    communicationnels)~;
  \item Les modalites d'intervention de chaque educateur specialise~;
  \item Les criteres d'evaluation des progres realises~;
  \item Le calendrier de revision et d'actualisation des objectifs~;
  \item Les recommandations adressees aux familles.
\end{itemize}
\end{encadre}

\section{Deroulement du stage}

Le stage s'est deroule du \textbf{12 mars au 12 juin 2026}, soit une
periode de trois mois au sein du centre de Martil.

\begin{table}[H]
\centering
\caption{Chronogramme des activites de stage}
\label{tab:chronogramme}
\begin{tabularx}{\textwidth}{llX}
\toprule
\textbf{Periode} & \textbf{Phase} & \textbf{Activites principales}\\
\midrule
12 mars -- 31 mars & Phase 1 & Integration, observation, audit de l'existant\\
1 avril -- 30 avril & Phase 2 & Conception et developpement de la plateforme\\
1 mai -- 25 mai & Phase 3 & Tests, deploiement, formation utilisateurs\\
26 mai -- 12 juin & Phase 4 & Evaluation de l'impact, redaction du rapport\\
\bottomrule
\end{tabularx}
\vspace{0.1cm}
\small\textit{Source~: Plan de projet elabore par l'auteur, 2026.}
\end{table}

%% ================================================================
\chapter{Innovation manageriale et transformation digitale~: cadre theorique}

\section{L'innovation manageriale}

\subsection{Definition}

L'innovation manageriale designe l'introduction de pratiques, processus
ou structures organisationnelles substantiellement nouveaux, dans le
but d'ameliorer la performance de l'organisation\footnote{Birkinshaw,
J., \& Mol, M. (2006). How management innovation happens. \textit{MIT
Sloan Management Review}, 47(4), 81--88.}.

\begin{table}[H]
\centering
\caption{Les quatre dimensions de l'innovation manageriale}
\label{tab:dimensions}
\begin{tabularx}{\textwidth}{lXX}
\toprule
\textbf{Dimension} & \textbf{Definition} & \textbf{Application El Awael}\\
\midrule
Processus & Nouveaux modes d'organisation du travail & Workflow numerique PEI\\
Structure & Nouvelles formes de coordination & Tableau de bord centralise\\
Pratiques & Nouvelles manieres de decider & KPIs et alertes automatiques\\
Systemes & Nouveaux outils de pilotage & Plateforme integree\\
\bottomrule
\end{tabularx}
\vspace{0.1cm}
\small\textit{Source~: Adapte de Birkinshaw \& Mol (2006), application auteur.}
\end{table}

\section{La transformation digitale}

\subsection{Definition et portee}

La transformation digitale est definie comme le processus par lequel
une organisation integre les technologies numeriques dans l'ensemble
de ses activites, entrainant des changements fondamentaux dans son
fonctionnement et dans la valeur qu'elle delivre\footnote{Westerman,
G., Bonnet, D., \& McAfee, A. (2014). \textit{Leading Digital}.
Harvard Business Press.}.

\begin{table}[H]
\centering
\caption{Trois niveaux du passage au numerique}
\label{tab:niveaux}
\begin{tabularx}{\textwidth}{lXX}
\toprule
\textbf{Niveau} & \textbf{Definition} & \textbf{Exemple}\\
\midrule
Numerisation & Conversion de supports physiques & Scanner des dossiers papier\\
Digitalisation & Usage du numerique pour ameliorer & Tableur Excel remplacant un registre\\
Transformation digitale & Refonte globale de l'organisation & Plateforme integree + KPIs\\
\bottomrule
\end{tabularx}
\vspace{0.1cm}
\small\textit{Source~: Adapte de Gartner (2021), application auteur.}
\end{table}

\section{Le Management 4.0}

Le Management 4.0 est une transposition des principes de l'Industrie 4.0
au champ du management organisationnel. Il se caracterise par l'integration
des technologies numeriques dans les processus de decision et de pilotage.

\begin{encadre}{Les cinq piliers du Management 4.0}
\begin{enumerate}
  \item \textbf{Pilotage par les donnees}~: toute decision s'appuie sur des
    indicateurs mesures en temps reel~;
  \item \textbf{Agilite organisationnelle}~: adaptation rapide aux besoins
    des beneficiaires~;
  \item \textbf{Collaboration numerique}~: partage fluide entre educateurs,
    administration et direction~;
  \item \textbf{Automatisation des taches a faible valeur ajoutee}~: liberer
    du temps pour l'accompagnement humain~;
  \item \textbf{Transparence et redevabilite}~: tracabilite complete des
    actions pour les partenaires et financeurs.
\end{enumerate}
\end{encadre}

\section{Specificites de la transformation digitale dans le secteur associatif}

La transformation digitale dans les associations a vocation sociale presente
des specificites importantes~:

\begin{itemize}
  \item \textbf{Contrainte budgetaire severe}~: toute solution doit etre
    gratuite ou a cout marginal nul~;
  \item \textbf{Resistance culturelle}~: les equipes peuvent percevoir le
    numerique comme une menace a la dimension humaine~;
  \item \textbf{Sensibilite des donnees}~: les dossiers des beneficiaires
    sont soumis aux obligations de protection les plus strictes~;
  \item \textbf{Exigence de reporting}~: les partenaires institutionnels
    exigent des rapports d'activite de plus en plus detailles.
\end{itemize}

%% ================================================================
\part{Diagnostic organisationnel et conception de la solution}

%% ================================================================
\chapter{Diagnostic de la situation existante}

\section{Methodologie du diagnostic}

Pour etablir un diagnostic organisationnel rigoureux, j'ai mobilise~:

\begin{itemize}
  \item \textbf{Observation participante}~: presence quotidienne pendant
    quatre semaines~;
  \item \textbf{Entretiens semi-directifs}~: avec la directrice, trois
    educateurs et le responsable administratif~;
  \item \textbf{Analyse documentaire}~: examen des registres papier et
    fichiers informatiques~;
  \item \textbf{Cartographie des processus}~: reconstitution du flux
    d'information de la gestion d'un beneficiaire.
\end{itemize}

\section{Dysfonctionnements critiques identifies}

\begin{table}[H]
\centering
\caption{Synthese des dysfonctionnements critiques}
\label{tab:dysfonctionnements}
\begin{tabularx}{\textwidth}{llXl}
\toprule
\textbf{N} & \textbf{Priorite} & \textbf{Dysfonctionnement} & \textbf{Impact}\\
\midrule
1 & \textcolor{rougeKO}{\bfseries CRITIQUE} & 44 mots de passe en clair dans le HTML & Securite\\
2 & \textcolor{rougeKO}{\bfseries CRITIQUE} & Donnees stockees localement (perte possible) & Integrite\\
3 & \textcolor{orange}{\bfseries IMPORTANT} & Absence d'indicateurs de pilotage (KPIs) & Decision\\
4 & \textcolor{orange}{\bfseries IMPORTANT} & Alertes uniquement reactives, non proactives & Prevention\\
5 & \textcolor{orange}{\bfseries IMPORTANT} & Export limite a l'impression & Reporting\\
6 & \textcolor{vertOK}{\bfseries MOYEN} & Interface non professionnelle (alertes navigateur) & Experience\\
\bottomrule
\end{tabularx}
\vspace{0.1cm}
\small\textit{Source~: Audit de terrain realise par l'auteur, mars 2026.}
\end{table}

\section{Analyse SWOT}

\begin{table}[H]
\centering
\caption{Analyse SWOT de l'Association El Awael --- Centre de Martil}
\label{tab:swot}
\begin{tabularx}{\textwidth}{XX}
\toprule
\textcolor{vertOK}{\bfseries Forces} & \textcolor{rougeKO}{\bfseries Faiblesses}\\
\midrule
\begin{itemize}[leftmargin=*,topsep=0pt]
  \item Equipe educative engagee et experimentee
  \item Reconnaissance institutionnelle locale
  \item Approche PEI individualisee
  \item Forte legitimite aupres des familles
\end{itemize}
&
\begin{itemize}[leftmargin=*,topsep=0pt]
  \item Gestion entierement manuelle
  \item Absence de SI centralise
  \item Securite des donnees insuffisante
  \item Reporting chronophage et peu fiable
\end{itemize}\\
\midrule
\textcolor{bleu}{\bfseries Opportunites} & \textcolor{orange}{\bfseries Menaces}\\
\midrule
\begin{itemize}[leftmargin=*,topsep=0pt]
  \item Outils cloud gratuits (Firebase)
  \item Appui academique (stage Master)
  \item Contexte national favorable a l'inclusion
\end{itemize}
&
\begin{itemize}[leftmargin=*,topsep=0pt]
  \item Risque de perte de donnees
  \item Turn-over du personnel
  \item Resistance au changement
\end{itemize}\\
\bottomrule
\end{tabularx}
\vspace{0.1cm}
\small\textit{Source~: Elaboration de l'auteur d'apres analyse de terrain, 2026.}
\end{table}

\section{Analyse des besoins (MoSCoW)}

\begin{table}[H]
\centering
\caption{Analyse MoSCoW des besoins}
\label{tab:moscow}
\begin{tabularx}{\textwidth}{lX}
\toprule
\textbf{Priorite} & \textbf{Besoins identifies}\\
\midrule
\textcolor{rougeKO}{\bfseries Must have} &
  Authentification securisee, dossiers beneficiaires en ligne,
  suivi presences, stockage cloud securise\\
\midrule
\textcolor{orange}{\bfseries Should have} &
  Tableau de bord KPIs, alertes automatiques, export PDF/Excel\\
\midrule
\textcolor{bleu}{\bfseries Could have} &
  Simulation strategique, heatmap des risques, journal d'audit,
  dashboard institutionnel\\
\midrule
\textcolor{grisTxt}{\bfseries Won't have} &
  Application mobile native, integration systemes ministeriels\\
\bottomrule
\end{tabularx}
\vspace{0.1cm}
\small\textit{Source~: Elaboration de l'auteur d'apres diagnostic, 2026.}
\end{table}

%% ================================================================
\chapter{Methodologie et conception de la solution digitale}

\section{Methodologie agile adaptee}

Face a la complexite des besoins, j'ai opte pour une methodologie
de developpement agile, decoupee en \textbf{quatre sprints}~:

\begin{table}[H]
\centering
\caption{Decoupage en sprints du projet}
\label{tab:sprints}
\begin{tabularx}{\textwidth}{llXl}
\toprule
\textbf{Sprint} & \textbf{Periode} & \textbf{Livrable} & \textbf{Validation}\\
\midrule
Sprint 1 & 1--14 avril & Authentification Firebase + login & Directrice\\
Sprint 2 & 15--28 avril & Gestion beneficiaires + PEI & Educateurs\\
Sprint 3 & 1--15 mai & KPIs + alertes + simulation & Direction\\
Sprint 4 & 16--25 mai & Heatmap + audit + deploiement & Equipe\\
\bottomrule
\end{tabularx}
\vspace{0.1cm}
\small\textit{Source~: Plan de projet elabore par l'auteur, 2026.}
\end{table}

\section{Architecture technique}

La plateforme repose sur une architecture \textbf{client-side} s'appuyant
sur des services cloud gratuits, eliminant le besoin d'un serveur dedie~:

\begin{figure}[H]
\centering
\begin{tcolorbox}[colback=bleuClair!20,colframe=bleu,boxrule=0.8pt,arc=3pt,width=0.9\textwidth]
\centering\small
\begin{tabular}{ccc}
\fbox{\parbox{4cm}{\centering\bfseries Navigateur Web\\(HTML5 + CSS3 + JS pur)}} &
$\longleftrightarrow$ &
\fbox{\parbox{4cm}{\centering\bfseries Firebase\\Firestore DB\\Authentication\\Storage}}\\[0.5cm]
$\uparrow$ & & \\[0.1cm]
\fbox{\parbox{4cm}{\centering\bfseries GitHub Pages\\Deploiement statique gratuit}} & &
\end{tabular}
\end{tcolorbox}
\caption{Architecture technique globale de la plateforme ElAwael}
\label{fig:architecture}
\small\textit{Source~: Elaboration de l'auteur, 2026.}
\end{figure}

\section{Choix technologiques}

\begin{table}[H]
\centering
\caption{Technologies retenues et justification}
\label{tab:techno}
\begin{tabularx}{\textwidth}{llXl}
\toprule
\textbf{Technologie} & \textbf{Role} & \textbf{Justification} & \textbf{Cout}\\
\midrule
Firebase Firestore & Base de donnees & Temps reel, scalable, securise & 0 MAD\\
Firebase Auth & Authentification & Sessions, roles, reset mdp par email & 0 MAD\\
Firebase Storage & Fichiers & Stockage documents beneficiaires & 0 MAD\\
GitHub Pages & Hebergement & Deploiement statique gratuit & 0 MAD\\
GitHub Actions & CI/CD & Deploiement automatique a chaque commit & 0 MAD\\
Chart.js 4.4 & Visualisation & Graphiques interactifs, bibliotheque legere & 0 MAD\\
JavaScript pur & Logique & Aucune dependance framework & 0 MAD\\
\bottomrule
\end{tabularx}
\vspace{0.1cm}
\small\textit{Source~: Analyse des options effectuee par l'auteur, 2026.}
\end{table}

\section{Systeme de roles et gestion des acces (RBAC)}

\begin{table}[H]
\centering
\caption{Matrice des permissions par role}
\label{tab:rbac}
\begin{tabularx}{\textwidth}{Xcccc}
\toprule
\textbf{Fonctionnalite} & \textbf{Admin} & \textbf{Gestionnaire} & \textbf{Educateur} & \textbf{Consultant}\\
\midrule
Voir les KPIs & \cmark & \cmark & \cmark & \cmark\\
Dossiers beneficiaires & \cmark & \cmark & \cmark & Masque\\
Modifier les donnees & \cmark & \cmark & \cmark & \xmark\\
Supprimer des donnees & \cmark & \xmark & \xmark & \xmark\\
Gerer les utilisateurs & \cmark & \xmark & \xmark & \xmark\\
Exporter PDF/Excel & \cmark & \cmark & \xmark & \xmark\\
Journal d'audit & \cmark & \cmark & \xmark & \xmark\\
\bottomrule
\end{tabularx}
\vspace{0.1cm}
\small\textit{Source~: Specifications fonctionnelles, auteur, 2026.}
\end{table}

%% ================================================================
\part{Realisation de la plateforme et evaluation de l'impact}

%% ================================================================
\chapter{La plateforme ElAwael Digital Platform}

\section{Vue d'ensemble}

La plateforme ElAwael Digital Platform est un \textbf{systeme d'information
integre}, accessible depuis n'importe quel navigateur web, sans installation
prealable, couvrant l'ensemble du cycle operationnel de l'association.

\begin{figure}[H]
\centering
\safeimg[0.92\textwidth]{screenshots/overview.png}{Vue d'ensemble de la plateforme}
\caption{Vue d'ensemble de la plateforme ElAwael Digital Platform}
\label{fig:overview}
\small\textit{Source~: Capture d'ecran realisee par l'auteur, mai 2026.}
\end{figure}

\section{Module d'authentification et securite}

\begin{figure}[H]
\centering
\safeimg[0.82\textwidth]{screenshots/sc_login.png}{Page de connexion securisee}
\caption{Page de connexion securisee}
\label{fig:login}
\small\textit{Source~: Capture d'ecran realisee par l'auteur, avril 2026.}
\end{figure}

\begin{encadre}{Couches de securite implementees}
\begin{enumerate}
  \item Authentification Firebase avec tokens JWT auto-renouvelables~;
  \item Persistance de session via IndexedDB (reste connecte apres fermeture)~;
  \item Regles Firestore~: acces zero pour les non-authentifies~;
  \item Protection anti-boucle de redirection~;
  \item Reset de mot de passe automatique par email.
\end{enumerate}
\end{encadre}

\vspace{0.3cm}
Avant l'implementation, la situation securitaire etait critique~:
\textbf{44 mots de passe en clair} etaient codes directement dans le
fichier HTML principal, accessibles a toute personne consultant le code source.

\section{Tableau de bord et indicateurs cles de performance}

\begin{figure}[H]
\centering
\safeimg[0.92\textwidth]{screenshots/sc_kpi.png}{Tableau de bord KPIs}
\caption{Tableau de bord decisionnel --- Indicateurs Cles de Performance}
\label{fig:kpi}
\small\textit{Source~: Capture d'ecran realisee par l'auteur, mai 2026.}
\end{figure}

\begin{figure}[H]
\centering
\safeimg[0.92\textwidth]{screenshots/dashboard2.png}{Vue complete tableau de bord}
\caption{Vue complete du tableau de bord administrateur}
\label{fig:dashboard2}
\small\textit{Source~: Capture d'ecran realisee par l'auteur, mai 2026.}
\end{figure}

\begin{table}[H]
\centering
\caption{Les huit KPIs principaux et leur formule de calcul}
\label{tab:kpis}
\begin{tabularx}{\textwidth}{lXl}
\toprule
\textbf{KPI} & \textbf{Formule} & \textbf{Seuil critique}\\
\midrule
Realisation PEI & Objectifs realises / Objectifs totaux $\times$ 100 & $<$ 50\%\\
Taux de presence & Presences / Seances prevues $\times$ 100 & $<$ 70\%\\
Beneficiaires a risque & Count(score\_risque $\geq$ 61) & $>$ 20\%\\
Ratio cadre/benef. & Nb beneficiaires / Nb cadres actifs & $>$ 8\\
Cadres actifs & Count(seances ce mois $>$ 0) & $<$ 80\%\\
Alertes non resolues & Count(alertes.statut = ouvert) & $>$ 5\\
Sans seance ce mois & Count(seances\_mois = 0) & $>$ 3\\
Dossiers complets & Dossiers complets / Total $\times$ 100 & $<$ 85\%\\
\bottomrule
\end{tabularx}
\vspace{0.1cm}
\small\textit{Source~: Specifications fonctionnelles, auteur, 2026.}
\end{table}

\subsection{Le Score Organisationnel Global (SOG)}

\begin{figure}[H]
\centering
\safeimg[0.85\textwidth]{screenshots/sc_sog.png}{Score Organisationnel Global}
\caption{Jauge du Score Organisationnel Global (SOG)}
\label{fig:sog}
\small\textit{Source~: Capture d'ecran realisee par l'auteur, mai 2026.}
\end{figure}

\begin{encadre}{Formule du Score Organisationnel Global (SOG)}
\[
SOG = \bigl(\bar{R}_{PEI} \times 40\%\bigr)
    + \bigl(\bar{P}_{presence} \times 30\%\bigr)
    + \bigl(\bar{I}_{indep.} \times 20\%\bigr)
    + \bigl(T_{dossiers} \times 10\%\bigr)
\]
\textbf{Interpretation~:}
\textcolor{vertOK}{$\geq 70$ = Excellent}~|~%
\textcolor{orange}{$\geq 40$ = Bien}~|~%
\textcolor{rougeKO}{$< 40$ = A ameliorer}
\end{encadre}

\section{Gestion des beneficiaires}

\begin{figure}[H]
\centering
\safeimg[0.92\textwidth]{screenshots/sc_benef.png}{Module beneficiaires}
\caption{Module de gestion des beneficiaires}
\label{fig:benef}
\small\textit{Source~: Capture d'ecran realisee par l'auteur, mai 2026.}
\end{figure}

\begin{figure}[H]
\centering
\safeimg[0.92\textwidth]{screenshots/beneficiaires.png}{Liste beneficiaires}
\caption{Liste complete avec indicateurs de risque individuel}
\label{fig:liste-benef}
\small\textit{Source~: Capture d'ecran realisee par l'auteur, mai 2026.}
\end{figure}

\begin{table}[H]
\centering
\caption{Criteres de calcul du score de risque individuel}
\label{tab:risque}
\begin{tabularx}{\textwidth}{Xc}
\toprule
\textbf{Critere} & \textbf{Points}\\
\midrule
Taux de presence $<$ 50\% & +40\\
Taux de presence entre 50\% et 70\% & +25\\
Taux de realisation PEI $<$ 30\% & +35\\
Taux de realisation PEI entre 30\% et 50\% & +20\\
Dossier incomplet (champs manquants) & +15\\
Aucune seance enregistree ce mois & +10\\
\midrule
\textbf{Interpretation} : 0--30 = Stable / 31--60 = Surveiller / 61--100 = Urgent & \\
\bottomrule
\end{tabularx}
\vspace{0.1cm}
\small\textit{Source~: Specifications fonctionnelles, auteur, 2026.}
\end{table}

\section{Module de simulation strategique}

\begin{figure}[H]
\centering
\safeimg[0.92\textwidth]{screenshots/sc_simulation.png}{Simulation strategique}
\caption{Module de simulation strategique --- leviers d'action et score simule}
\label{fig:simulation}
\small\textit{Source~: Capture d'ecran realisee par l'auteur, mai 2026.}
\end{figure}

Le simulateur offre trois leviers reglables~: amelioration de la presence
(0 a +30\%), de la realisation PEI (0 a +30\%) et de l'autonomie ABLLS
(0 a +20\%). Le score simule est recalcule en temps reel.

\section{Heatmap des risques et alertes}

\begin{figure}[H]
\centering
\safeimg[0.92\textwidth]{screenshots/alertes.png}{Heatmap et alertes}
\caption{Heatmap des risques et centre d'alertes --- Vue par classe}
\label{fig:alertes}
\small\textit{Source~: Capture d'ecran realisee par l'auteur, mai 2026.}
\end{figure}

\section{Journal d'audit}

\begin{table}[H]
\centering
\caption{Typologies d'evenements traces dans le journal d'audit}
\label{tab:audit}
\begin{tabularx}{\textwidth}{lX}
\toprule
\textbf{Categorie} & \textbf{Exemples d'evenements traces}\\
\midrule
Connexion & Connexion reussie, echec d'authentification, deconnexion\\
Beneficiaires & Creation, modification, suppression d'un dossier\\
PEI & Mise a jour des objectifs, evaluation saisie\\
Presences & Saisie de presence, modification d'une entree\\
Documents & Upload, telechargement, suppression de documents\\
Exports & Export Excel, export PDF, rapport institutionnel\\
Suppressions & Toute suppression avec identification de l'auteur\\
\bottomrule
\end{tabularx}
\vspace{0.1cm}
\small\textit{Source~: Specifications du module d'audit, auteur, 2026.}
\end{table}

%% ================================================================
\chapter{Impact organisationnel et humain}

\section{Evaluation de la performance}

\begin{table}[H]
\centering
\caption{Comparaison des indicateurs avant et apres deploiement}
\label{tab:impact}
\begin{tabularx}{\textwidth}{Xccl}
\toprule
\textbf{Indicateur} & \textbf{Avant} & \textbf{Apres} & \textbf{Evolution}\\
\midrule
Production d'un rapport direction & 3--4 heures & $<$ 5 min & \textcolor{vertOK}{$\downarrow$ 97\%}\\
Completude des dossiers & $\approx$ 45\% & $>$ 92\% & \textcolor{vertOK}{$\uparrow$ +47 pts}\\
Detection d'un beneficiaire a risque & Plusieurs semaines & Temps reel & \textcolor{vertOK}{Immediat}\\
Tracabilite des actions & Nulle & Complete & \textcolor{vertOK}{100\%}\\
Securite des donnees & Insuffisante & Conforme & \textcolor{vertOK}{OK}\\
Accessibilite (deplacement) & Locale uniquement & Tout navigateur & \textcolor{vertOK}{Cloud}\\
Cout de la solution & -- & 0 MAD/mois & \textcolor{vertOK}{Gratuit}\\
\bottomrule
\end{tabularx}
\vspace{0.1cm}
\small\textit{Source~: Mesures realisees par l'auteur, mars et juin 2026.}
\end{table}

\section{Gain de temps pour les equipes}

\begin{table}[H]
\centering
\caption{Gain de temps hebdomadaire par poste}
\label{tab:temps}
\begin{tabularx}{\textwidth}{Xcc}
\toprule
\textbf{Tache automatisee} & \textbf{Avant} & \textbf{Apres}\\
\midrule
Saisie des presences & 45 min/sem. & 10 min/sem.\\
Mise a jour des objectifs PEI & 2h/sem. & 30 min/sem.\\
Preparation des bilans & 3h/mois & 20 min/mois\\
Reporting direction & 4h/mois & 5 min/mois\\
Recherche dans les dossiers & 20 min/recherche & $<$ 1 min\\
\midrule
\textbf{Total estime} & \textbf{7h/semaine/poste} & \textbf{1h30/semaine/poste}\\
\bottomrule
\end{tabularx}
\vspace{0.1cm}
\small\textit{Source~: Entretiens de retour d'experience, juin 2026.}
\end{table}

\section{Retour d'experience des utilisateurs}

\begin{table}[H]
\centering
\caption{Resultats du questionnaire de satisfaction (n=5)}
\label{tab:satisfaction}
\begin{tabularx}{\textwidth}{Xccc}
\toprule
\textbf{Critere evalue} & \textbf{Insatisfait} & \textbf{Neutre} & \textbf{Satisfait}\\
\midrule
Facilite de prise en main & 0\% & 20\% & 80\%\\
Clarte de l'interface & 0\% & 0\% & 100\%\\
Utilite des KPIs & 0\% & 0\% & 100\%\\
Amelioration du suivi & 0\% & 20\% & 80\%\\
Gain de temps percu & 0\% & 0\% & 100\%\\
Souhait de continuer & 0\% & 0\% & 100\%\\
\bottomrule
\end{tabularx}
\vspace{0.1cm}
\small\textit{Source~: Questionnaire administre par l'auteur, juin 2026~; n=5.}
\end{table}

\begin{boitecitation}
\itshape
\og{}Avant, je savais intuitivement quels enfants avaient besoin de plus
d'attention. Maintenant, j'ai les chiffres qui le confirment, et je peux
les montrer a nos partenaires.\fg{}\\
\raggedleft --- Mme Kaoutar Dghoughi, directrice, entretien du 5 juin 2026.
\end{boitecitation}

\section{Recommandations strategiques}

\begin{enumerate}
  \item \textbf{Designer un referent numerique interne} charge de la
    maintenance et de la formation des nouveaux arrivants~;
  \item \textbf{Systematiser la saisie} via une procedure operationnelle
    standard definissant qui saisit quoi, quand et comment~;
  \item \textbf{Etendre aux autres centres} (M'diq et Fnideq) avec des
    adaptations mineures~;
  \item \textbf{Envisager une version mobile} (Progressive Web App) pour
    la saisie des presences sur smartphone~;
  \item \textbf{Formaliser le reporting institutionnel} en exploitant le
    dashboard financeur pour les partenaires.
\end{enumerate}

%% ================================================================
\chapter*{Conclusion generale}
\addcontentsline{toc}{chapter}{Conclusion generale}
\markboth{Conclusion generale}{Conclusion generale}

Ce stage de fin d'etudes au sein de l'Association El Awael --- centre de
Martil --- a constitue une experience fondatrice a plus d'un titre. Il a
permis d'eprouver in vivo la puissance et les limites de la transformation
digitale lorsqu'elle est conduite dans un contexte associatif caracterise
par des ressources limitees, mais une mission sociale d'une haute noblesse.

\vspace{0.4cm}
La question centrale posee en introduction a trouve des elements de
reponse concrets. La plateforme ElAwael Digital Platform demontre qu'il
est possible, avec des outils entierement gratuits et un investissement en
competences plutot qu'en capital financier, de passer d'une gestion
intuitive et papier a un pilotage base sur les donnees, sans sacrifier la
dimension humaine qui constitue l'identite profonde de l'association.

\vspace{0.4cm}
Les trois apports principaux de ce travail sont~:

\begin{enumerate}
  \item \textbf{Un apport theorique}~: la demonstration que les concepts
    du Management 4.0 sont transposables au secteur associatif avec des
    adaptations appropriees~;

  \item \textbf{Un apport methodologique}~: l'articulation d'un diagnostic
    rigoureux avec un developpement agile centre sur les utilisateurs reels~;

  \item \textbf{Un apport pratique}~: une plateforme de gestion numerique
    integree, operationnelle, securisee, bilingue et accessible, couvrant
    l'ensemble du cycle de vie organisationnel de l'Association El Awael.
\end{enumerate}

\vspace{0.4cm}
Si ce stage a confirme une conviction, c'est celle-ci~: la technologie
n'est pas une fin en soi, mais un \textbf{moyen au service d'une mission}.
Et cette mission --- accompagner avec dignite des enfants en situation de
handicap vers le meilleur developpement possible --- merite tous les outils
que l'ingenuosite humaine peut mobiliser a son service.

%% ================================================================
\begin{thebibliography}{99}
\addcontentsline{toc}{chapter}{Bibliographie}

\bibitem{birkinshaw2006}
Birkinshaw, J., \& Mol, M. (2006). How management innovation happens.
\textit{MIT Sloan Management Review}, 47(4), 81--88.

\bibitem{hamel2006}
Hamel, G. (2006). The why, what, and how of management innovation.
\textit{Harvard Business Review}, 84(2), 72--84.

\bibitem{westerman2014}
Westerman, G., Bonnet, D., \& McAfee, A. (2014).
\textit{Leading Digital: Turning Technology into Business Transformation}.
Boston~: Harvard Business Press.

\bibitem{roblek2016}
Roblek, V., Mesko, M., \& Krapez, A. (2016).
A Complex View of Industry 4.0. \textit{SAGE Open}, 6(2), 1--11.

\bibitem{bentaleb2019}
Bentaleb, A., Benali, N., \& Hmioui, A. (2019).
Enjeux de la transformation digitale pour le management des organisations.
\textit{Revue Internationale des Sciences de Gestion}, 2(3), 114--132.

\bibitem{drucker1985}
Drucker, P. (1985).
\textit{Innovation and Entrepreneurship}. New York~: Harper \& Row.

\bibitem{gartner2021}
Gartner. (2021). \textit{Gartner Glossary: Digitalization}.
Stamford~: Gartner Inc.

\bibitem{onu2006}
ONU. (2006). \textit{Convention relative aux droits des personnes
handicapees}. New York~: ONU. [Ratifiee par le Maroc en 2009]

\bibitem{sgg2024}
Secretariat General du Gouvernement. (2024).
\textit{Rapport annuel sur le mouvement associatif au Maroc}.
Rabat~: Imprimerie Officielle.

\bibitem{firebase2024}
Google Firebase. (2024). \textit{Firebase Documentation}.
\url{https://firebase.google.com/docs}

\end{thebibliography}

%% ================================================================
%%  ANNEXES
%% ================================================================
\appendix
\chapter{Captures d'ecran complementaires}

\begin{figure}[H]
\centering
\safeimg[0.92\textwidth]{screenshots/portal.png}{Portail principal}
\caption{Vue portail de la plateforme ElAwael}
\label{fig:portal}
\small\textit{Source~: Capture d'ecran realisee par l'auteur, mai 2026.}
\end{figure}

\begin{figure}[H]
\centering
\safeimg[0.92\textwidth]{screenshots/dashboard.png}{Dashboard educateur}
\caption{Tableau de bord educateur}
\label{fig:dash-educ}
\small\textit{Source~: Capture d'ecran realisee par l'auteur, mai 2026.}
\end{figure}

\chapter{Questionnaire de retour d'experience}

\textbf{1. Satisfaction globale vis-a-vis de la plateforme~?}\\
$\square$ Tres insatisfait \quad $\square$ Insatisfait \quad $\square$ Neutre \quad
$\square$ Satisfait \quad $\square$ Tres satisfait

\vspace{0.4cm}
\textbf{2. La plateforme a-t-elle ameliore votre efficacite~?}\\
$\square$ Pas du tout \quad $\square$ Peu \quad $\square$ Moyennement \quad
$\square$ Beaucoup \quad $\square$ Tres fortement

\vspace{0.4cm}
\textbf{3. Gain de temps hebdomadaire estime~?}\\
$\square$ Moins d'1h \quad $\square$ 1--3h \quad $\square$ 3--5h \quad $\square$ Plus de 5h

\vspace{0.4cm}
\textbf{4. L'interface est-elle intuitive sans formation longue~?}\\
$\square$ Oui, totalement \quad $\square$ Avec quelques difficultes \quad
$\square$ Formation necessaire

\vspace{0.4cm}
\textbf{5. Souhaitez-vous que la plateforme soit etendue aux autres centres~?}\\
$\square$ Oui, certainement \quad $\square$ Plutot oui \quad $\square$ Neutre \quad $\square$ Non

\vspace{0.4cm}
\textbf{6. Suggestions~:}\\
\rule{\textwidth}{0.4pt}\\[0.3cm]
\rule{\textwidth}{0.4pt}

\chapter{Regles de securite Firestore}

\begin{verbatim}
rules_version = '2';
service cloud.firestore {
  match /databases/{database}/documents {
    match /beneficiaires/{docId} {
      allow read: if request.auth != null;
      allow write: if request.auth != null
        && get(/databases/$(database)/documents/
               users/$(request.auth.uid)).data.role
           in ['ADMIN', 'GESTIONNAIRE', 'EDUCATEUR'];
    }
    match /audit_logs/{logId} {
      allow read: if request.auth != null;
      allow write: if request.auth != null;
    }
    match /users/{userId} {
      allow read: if request.auth != null;
      allow write: if get(/databases/$(database)/documents/
        users/$(request.auth.uid)).data.role == 'ADMIN';
    }
  }
}
\end{verbatim}

\end{document}
